\bookchapter{第六十五章：プロジェクト1——わたしの家計簿（ローカル版）}

\chapterintro{あなたの初めての「まともなプロジェクト」へようこそ。この家計簿はネットにはつながりませんが、中身は一通り揃っています：1件の記録ができ、履歴が見られて、月ごとに集計できて、削除·編集もできます。完成すれば、「データ ロジック 画面」の全工程をひととおり歩いたことになります。コードをいきなり写すのは急がないで——まずよく考えてから手を動かす。それが私たちのいつもの原則です。}

\booksection{要件を考える：まず5つの質問}

\begin{center}
\begin{tabularx}{\linewidth}{|X|X|}
\hline
\textbf{要素} & \textbf{答え} \tabularnewline
\hline
Who（誰が使うか） & 自分用。家族や友だちでもOK \tabularnewline
What（何を作るか） & 収支を記録、履歴を見る、月ごとにまとめる、カテゴリ別に集計 \tabularnewline
Why（なぜ作るか） & 練習＋お金を何に使ったかをはっきりさせる \tabularnewline
How（どう作るか） & Swift＋SwiftUI＋SwiftData＋MVVM。完全ローカル、サーバーなし \tabularnewline
Constraint（制約） & 1週間で完成。iOS 17以上をサポート（SwiftDataの要件） \tabularnewline
\hline
\end{tabularx}
\end{center}

\begin{center}\small\color{gray}家計簿の要件5要素\end{center}

そして「どの段階でも使える」小さなステップに分解します：(1)SwiftDataのデータ層を作る（1件保存できる）。(2)リストページ（履歴が見られる）。(3)入力ポップアップ（記録を追加できる）。(4)削除と編集。(5)月次集計カード。1ステップ終えるごとに、アプリは「使えるもの」になります。

\booksection{データ層：記録に「家」を用意する}

まず1件の記録がどんな形かを定義します——フィールドは十分に、でも欲張らない：

\begin{lstlisting}[language=Swift]
import SwiftData

@Model
class Record {
    var type: String      // "income" 収入 | "expense" 支出
    var amount: Double    // 金額（元）
    var category: String  // カテゴリー：食費、交通、給料……
    var note: String      // メモ（任意）
    var date: Date        // 記帳日

    init(type: String, amount: Double, category: String, note: String = "", date: Date = Date()) {
        self.type = type
        self.amount = amount
        self.category = category
        self.note = note
        self.date = date
    }
}
\end{lstlisting}

\begin{tcbinfo}{ワンポイント：金額はなぜDouble？}
本格的な金融アプリは、浮動小数点の誤差（0.1＋0.2$\neq$0.3という昔からの問題）を避けるためにDecimal（正確な十進数）を使います。練習用プロジェクトならDoubleで十分です。いつか本物の金融商品を作るときは、Decimalか「分単位のInt」に切り替えるのを忘れずに——これは先人たちが何度も踏んできた落とし穴です。

\end{tcbinfo}

\booksection{MVVM：画面とロジックの「世帯分け」}

画面にデータベースを直接操作させる？　動きはしますが、画面のコードが長ったらしくなります。MVVMアーキテクチャなら、それぞれが持ち場を守ります：View（画面）は描くことだけ、ViewModel（ロジックの中枢）はデータを担当、Model（SwiftData）は保存を担当。

View（SwiftUIの画面）

ViewModel（@Observableのロジック中枢）

Model（SwiftDataのデータ層）

\begin{center}\small\color{gray}図4：MVVMの3層構造——データは一方向に流れ、各層は互いに踏み込まない\end{center}

\begin{lstlisting}[language=Swift]
import Observation

@Observable
class RecordViewModel {
    var records: [Record] = []
    private let context: ModelContext

    init(context: ModelContext) { self.context = context }

    func load() {
        // SwiftDataの#Predicateで絞り込み + 並べ替え
        let descriptor = FetchDescriptor<Record>(
            sortBy: [SortDescriptor(\.date, order: .reverse)]
        )
        records = (try? context.fetch(descriptor)) ?? []
    }

    func add(type: String, amount: Double, category: String, note: String) {
        let r = Record(type: type, amount: amount, category: category, note: note)
        context.insert(r)
        try? context.save()
        load()                       // 読み込み直す。画面が更新される
    }

    func delete(_ record: Record) {
        context.delete(record)
        try? context.save()
        load()
    }
}
\end{lstlisting}

画面の中で@State private var vm = RecordViewModel(...)としておき、vm.recordsをそのままリストに描画——データが変わると、画面も自動で更新されます。

\booksection{画面：リストページ＋入力ポップアップ}

\begin{lstlisting}[language=Swift]
struct RecordListView: View {
    @State private var vm: RecordViewModel
    @State private var showAdd = false

    var body: some View {
        NavigationStack {
            List {
                Section("今月の残高") {          // 集計カード
                    SummaryRow(records: vm.records)
                }
                Section("明細") {
                    ForEach(vm.records) { r in
                        RecordRow(record: r)
                    }
                    .onDelete { idx in
                        idx.forEach { vm.delete(vm.records[$0]) }
                    }
                }
            }
            .navigationTitle("私の家計簿")
            .toolbar {
                Button { showAdd = true } label: {
                    Image(systemName: "plus")
                }
            }
            .sheet(isPresented: $showAdd) {
                AddRecordView(vm: vm)          // 入力ダイアログ
            }
        }
    }
}
\end{lstlisting}

入力ポップアップのポイント：金額入力欄のキーボードは.keyboardType(.decimalPad)（数字キーボード）。収入／支出はセグメントコントロール（Picker）で、カテゴリはメニューで選択。「保存」を押したらvm.add(...)を呼び、dismiss()で閉じます。

\begin{tcbwarn}{注意：入力チェックの手抜きは禁物}
金額入力欄は必ずチェックしましょう：空の値、負の数、“abc”のような不正な入力はすべてはじく（Double(text)で判定し、失敗するとnilが返ります）。アプリがクラッシュするいちばんよくある原因のひとつが、「ユーザーの入力をお行儀のいい子だと思い込む」ことです。忘れずに：ユーザーが入力したものは、すべて敵として審査する。

\end{tcbwarn}

\booksection{編集、削除、集計}

\begin{itemize}
  \item 削除：Listの.onDeleteで左スワイプ削除＋確認ポップアップ（うっかり削除は、家計簿系アプリのユーザーがいちばん恨むものです）。
  \item 編集：行をタップして編集ポップアップへ。元の値を先に入れておく（1文字直すほうが、打ち直すよりずっとラク）。
  \item 集計：filter＋reduceで「今月の収入／支出／残高」を計算して、上部のカードに表示します：
\end{itemize}

\begin{lstlisting}[language=Swift]
func monthSummary(_ records: [Record]) -> (income: Double, expense: Double) {
    let cal = Calendar.current
    let monthRecords = records.filter {
        cal.isDate($0.date, equalTo: Date(), toGranularity: .month)
    }
    let income = monthRecords
        .filter { $0.type == "income" }
        .reduce(0) { $0 + $1.amount }
    let expense = monthRecords
        .filter { $0.type == "expense" }
        .reduce(0) { $0 + $1.amount }
    return (income, expense)
}
\end{lstlisting}

\booksection{検収：3層チェック＋2つの上級チャレンジ}

プロジェクトが完成したら、「機能 境界 アーキテクチャ」の3層で検収します：記録できるか、削除できるか、アプリを再起動してもデータが残っているか（機能）。金額に0や負の数を入れてもクラッシュしないか（境界）。画面とデータが混ざっていないか、フィールドがハードコードされていないか（アーキテクチャ）。

基本版の検収をパスしたら、自分に2品の「上級メニュー」を出しましょう：

\begin{enumerate}
  \item カテゴリ別集計：「統計」タブを追加して、Swift Chartsでカテゴリ別支出を円グラフ／棒グラフに——Apple公式のグラフライブラリで、数行のコードでグラフが描けます。
  \item iCloud同期：SwiftDataのiCloud同期は、設定をいくつか入れるだけ——帳簿がiPhone／iPadの間で自動同期されて、一気に「高級感」が出ます。
\end{enumerate}

\begin{tcbpractice}{やってみよう：今週中に完成させる}
「データ層 ViewModel リストページ 入力ポップアップ 集計カード」の順に手を動かし、1日1ステップ。つまづいたら、エラーをそのまま検索エンジンやAIに投げましょう（第7編『AIプログラミング実践編』を覚えていますか？　問題を説明するときは「期待した動き、実際の動き、再現手順」をはっきりと）。完成したら、1週間これで家計簿をつけてみてください——自分で作ったアプリを使うと、達成感が次へ進む力になります。

\end{tcbpractice}

\begin{tcbsummary}{本章のまとめ}
\begin{itemize}
  \item 要件5要素 「どの段階でも使える」小ステップに分解してから、着手する
  \item データ層：@Modelモデル＋modelContainerコンテナ＋FetchDescriptorクエリ
  \item MVVMの分担：Viewは描くだけ、@ObservableなViewModelがロジック、SwiftDataが保存を担当
  \item 画面：List＋Section＋toolbar＋sheetの入力ポップアップ
  \item 3層検収（機能／境界／アーキテクチャ）＋上級チャレンジ（Swift Charts、iCloud同期）
\end{itemize}

\end{tcbsummary}
